% foundations/notation.tex

\chapter{Notation}
\label{chap:notation}

This chapter collects the notation used repeatedly throughout the
thesis. It introduces symbols and writing conventions only; the
engineering concepts represented by these symbols are developed in their
respective chapters.

% ============================================================
\section{Typographic Conventions}
% ============================================================

Scalar quantities are written in italic type, for example \(s\),
\(\theta\), and \(\kappa\). Vectors are written in bold type, for
example \(\mathbf t\), \(\mathbf n\), and \(\mathbf b\). Sets and
spaces are denoted by capital or calligraphic letters, while operators
are denoted by calligraphic letters whenever this improves the
conceptual distinction between an object and an operation.

A prime denotes differentiation with respect to arc-length unless
another independent variable is stated explicitly:
\[
f'(s)=\frac{\dd f}{\dd s}(s).
\]

% ============================================================
\section{Domains and Spaces}
% ============================================================

The arc-length domain of an alignment is an interval
\[
I\subseteq\mathbb{R}.
\]
The symbol
\[
\mathcal X
\]
denotes a state space appropriate to the context, while
\[
\mathcal Y
\]
denotes a space of evaluated or derived quantities. The space of
admissible curvature functions is denoted by
\[
\mathcal K.
\]

The intrinsic track-related coordinate domain is written as
\[
\Omega_{\mathrm{track}},
\]
and a realized world-space domain as
\[
\Omega_{\mathrm{world}}.
\]
The precise metric structure of these domains is introduced when metric
realization is developed.

% ============================================================
\section{Coordinates}
% ============================================================

The symbol
\[
s
\]
denotes intrinsic arc-length along the alignment. Track-related
coordinates are written as
\[
(s,d,h),
\]
where \(d\) is a lateral offset and \(h\) a vertical offset relative to
the applicable local track frame.

A point in world space is denoted by
\[
x\in\Omega_{\mathrm{world}},
\]
or by \(x_{\mathrm w}\) where an explicit distinction from intrinsic or
track-related coordinates is required.

% ============================================================
\section{Geometry Symbols}
% ============================================================

The planar alignment curve, heading, curvature, and cant rotation are
denoted by
\[
\gamma(s),
\qquad
\theta(s),
\qquad
\kappa(s),
\qquad
\phi(s).
\]
Cant-rate is written as
\[
\omega(s)=\phi'(s).
\]

The local tangent, lateral, and binormal directions are denoted by
\[
\mathbf t(s),
\qquad
\mathbf n(s),
\qquad
\mathbf b(s),
\]
and the corresponding local frame by
\[
F(s)=\bigl(\mathbf t(s),\mathbf n(s),\mathbf b(s)\bigr).
\]

A generic state is written as \(x(s)\), while a control is written as
\(u(s)\). The symbols \(\mathrm{pose2}(s)\) and
\(\mathrm{pose3}(s)\) refer to the planar and spatial railway-state
concepts defined in \cref{chap:state-model} and developed further in the
state part of the thesis.

% ============================================================
\section{Qualified Quantities}
% ============================================================

Subscripts qualify the engineering status or origin of a quantity. The
principal conventions are
\[
(\cdot)_{\mathrm{target}},
\qquad
(\cdot)_{\mathrm{real}},
\qquad
(\cdot)_{\mathrm{meas}}.
\]
Additional subscripts may identify a coordinate domain, frame, source,
or realization context. Such subscripts qualify a quantity; they do not
by themselves establish identity between different engineering
concepts.

% ============================================================
\section{Operator Notation}
% ============================================================

Track-to-world and world-to-track mappings are denoted by
\[
\mathcal P_{\mathrm{tw}}
\qquad\text{and}\qquad
\mathcal P_{\mathrm{wt}},
\]
respectively. Coordinate transformations between world-space
representations are denoted by
\[
\mathcal C.
\]

The symbols
\[
\mathcal E
\qquad\text{and}\qquad
\mathcal F
\]
denote evaluation and frame operators where no more specific operator
name is required.

Sampling, interpolation, mechanical response, local adjustment, and
realization are denoted by
\[
\mathcal S,
\qquad
\mathcal I,
\qquad
\mathcal M,
\qquad
\mathcal A,
\qquad
\mathcal R.
\]
Their domains, codomains, and engineering meaning are defined in the
chapters in which the corresponding operations are introduced.

% ============================================================
\section{Earth-related Symbols}
% ============================================================

Geodesic curvature and normal curvature are denoted by
\[
\kappa_g
\qquad\text{and}\qquad
\kappa_n.
\]
The ellipsoid surface normal is denoted by
\[
\mathbf n_{\mathrm E}.
\]
These symbols are used only after the applicable surface and realization
context have been established.
